\section{Guidance for Concrete Instantiations}\label{sec:concrete}

The $\UG$ construction places minimal requirements on its component
algorithms, allowing for diverse instantiations. Here we discuss some
possible concrete choices that satisfy the security assumptions of our
theorems, especially when targeting approximately 128-bit security.

\paragraph{Nominal Groups.}
For the nominal group $\nominalgroup$, both X25519~\cite{RFC7748} and the
NIST P-256 curve~\cite{NIST.SP.800.186} are suitable choices. It is generally
held that both curves achieve \sdh under standard assumptions, with the best
known classical attacks being generic discrete logarithm algorithms. The
nominal group abstraction from~\cite{EC:ABHKLR21} captures both Montgomery
curves like Curve25519 and short Weierstrass curves like P-256.

For X25519, the set of honest exponents $\honestexponent$ consists of 32-byte
strings with the standard clamping applied, while $\exponent =
\mathbb{Z}/p\mathbb{Z}$ for the curve's prime order subgroup. The statistical
distance $\Delta_{\nominalgroup}$ between honestly-generated and uniform
exponents is negligible ($< 2^{-128}$), as analyzed
in~\cite{EC:ABHKLR21,CiC:BCDKSV24}. Similar reasoning applies to P-256 when
using rejection sampling or other standard techniques to generate uniform
scalars: the two distributions $D_H$ and $D_U$ are identical, thus the
$\Delta_{\nominalgroup}$ between honestly-generated and uniform exponents is
0, and all elements have the same probability.

\paragraph{Post-Quantum KEMs.}
ML-KEM-768~\cite{FIPS203} is a popular choice for the \indcca \kem
component. NIST's standardization process included extensive analysis of
ML-KEM's \indcca security, reducing \indcca security to the hardness of
Module-LWE~\cite{EUROSP:BDKLLSSSS18}. ML-KEM-768 targets NIST Security Level
3, which should be approximately as difficult and expensive to attack as
running a key search on a 192-bit block cipher. The level 3 parameter set has
been paired with ~128-bit curves in other hybrid constructions such as X-Wing
to allow for some cushion and hedging against improvements in lattice
cryptanalysis in the long term. The correctness error $\delta$ is bounded by
$2^{-164}$, which is negligible in the security parameter.

\paragraph{Key Derivation Functions.}
For the hash function $H$ used as the KDF, both SHA3-256 (or SHAKE256, its
XOF form) and HKDF-SHA256 are appropriate choices.

\subparagraph{SHA-3/SHAKE256.}  The Keccak sponge construction underlying
SHA-3~\cite{FIPS202} provides a clean instantiation when modeling $H$ as a
random oracle. When keyed, SHAKE256 also achieves \prf security under the
assumption that the Keccak permutation behaves as an ideal
permutation~\cite{EC:BDPV08}.  Sponge-based constructions like SHA-3 have
been shown to be indifferentiable against classical~\cite{EC:BDPV08} as well
as quantum adversaries~\cite{EPRINT:ACMT25}.

\subparagraph{HKDF-SHA256.} In the standard model, HKDF~\cite{C:Krawczyk10}
instantiated with HMAC-SHA256 provides \prf security under the assumption
that HMAC-SHA256 is a secure \prf when keyed with a high-entropy
key. Krawczyk's analysis~\cite{C:Krawczyk10} establishes that HKDF-Expand is
a \prf when the input key material has sufficient min-entropy, which is
guaranteed in $\UG$ since $k_1$ is either an honestly-generated \kem shared
secret or uniform random. HKDF has been shown to be indifferentiable from a
random oracle under specific constraints~\cite{EUROSP:LipBlaBha19}: that HMAC
is indifferentiable from a random oracle, which for HMAC-SHA-256 has been
shown in ~\cite{C:DRST12}, assuming the compression function underlying
SHA-256 is a random oracle, which it is indifferentiably when used
prefix-free, and the values of HKDF's \texttt{IKM} input do not collide with
values of \texttt{info || 0x01}.


The choice of the KDF $H$'s security level ought to be made based on the
security level provided by the constituent \kem.
